% Shared colour palette for the ETHOS.TSAM manuscript figures.
%
% ── WHY THIS FILE EXISTS ──────────────────────────────────────────────────
% Reviewer request on the submitted manuscript:
%
%   "I would use as few different colors across your figures as possible. […]
%    I would make sure that the Jülich blue (RGB 2, 61, 107) is the leading
%    theme. That + a gray tone are the basic colors in all of your figures,
%    for shaded boxes e.g. at 70% transparency or so."
%
% Each figure used to declare its own \definecolor block, and they drifted:
% Fig. 1 carried two unrelated light blues plus a teal and an amber, Fig. 2 an
% amber and a warm amber tint. Declaring the palette ONCE and \input-ing it is
% what makes "the same colours in every figure" a property of the source
% rather than a thing to re-check by eye after every edit.
%
% ── THE PALETTE ───────────────────────────────────────────────────────────
% TWO base colours, and nothing else. Every fill in the figures is one of them
% mixed with white, which is exactly what "70 % transparency over a white
% page" means and what xcolor's `!nn` suffix computes: fzjbluetint is Jülich
% blue at 30 % opacity.
%
%   fzjblue  #023D6B  Jülich blue, RGB 2/61/107 — the leading theme. It draws
%                     the core: the tsam package itself, its internal edges,
%                     all primary text, and every mandatory feature.
%   fzjgrey  #3F4A55  the one grey — the periphery: external dependencies,
%                     optional or conditional elements, and legend chrome.
%
% ── THE CONSTRAINT THIS PUTS ON THE FIGURES ───────────────────────────────
% Both figures distinguish more than two things, and with only two hues the
% rest has to be carried by SHAPE, LINE STYLE or TYPOGRAPHY. That is a feature,
% not a compromise: all three survive a greyscale print, which a third hue
% would not. Concretely, what used to be colour is now
%
%   Fig. 1  object flow vs. dependency   solid vs. dashed (UML says so anyway)
%           pipeline phase vs. component round vs. sharp corners (UML Action)
%           optional stage               italic
%   Fig. 2  mandatory vs. optional       filled vs. open circle (FODA says so
%                                        anyway), reinforced by blue vs. grey
%
% If a figure ever needs a third colour, the honest fix is to move that
% distinction onto one of those channels — not to add the colour here.
%
% ── USAGE ─────────────────────────────────────────────────────────────────
% % Shared colour palette for the ETHOS.TSAM manuscript figures.
%
% ── WHY THIS FILE EXISTS ──────────────────────────────────────────────────
% Reviewer request on the submitted manuscript:
%
%   "I would use as few different colors across your figures as possible. […]
%    I would make sure that the Jülich blue (RGB 2, 61, 107) is the leading
%    theme. That + a gray tone are the basic colors in all of your figures,
%    for shaded boxes e.g. at 70% transparency or so."
%
% Each figure used to declare its own \definecolor block, and they drifted:
% Fig. 1 carried two unrelated light blues plus a teal and an amber, Fig. 2 an
% amber and a warm amber tint. Declaring the palette ONCE and \input-ing it is
% what makes "the same colours in every figure" a property of the source
% rather than a thing to re-check by eye after every edit.
%
% ── THE PALETTE ───────────────────────────────────────────────────────────
% TWO base colours, and nothing else. Every fill in the figures is one of them
% mixed with white, which is exactly what "70 % transparency over a white
% page" means and what xcolor's `!nn` suffix computes: fzjbluetint is Jülich
% blue at 30 % opacity.
%
%   fzjblue  #023D6B  Jülich blue, RGB 2/61/107 — the leading theme. It draws
%                     the core: the tsam package itself, its internal edges,
%                     all primary text, and every mandatory feature.
%   fzjgrey  #3F4A55  the one grey — the periphery: external dependencies,
%                     optional or conditional elements, and legend chrome.
%
% ── THE CONSTRAINT THIS PUTS ON THE FIGURES ───────────────────────────────
% Both figures distinguish more than two things, and with only two hues the
% rest has to be carried by SHAPE, LINE STYLE or TYPOGRAPHY. That is a feature,
% not a compromise: all three survive a greyscale print, which a third hue
% would not. Concretely, what used to be colour is now
%
%   Fig. 1  object flow vs. dependency   solid vs. dashed (UML says so anyway)
%           pipeline phase vs. component round vs. sharp corners (UML Action)
%           optional stage               italic
%   Fig. 2  mandatory vs. optional       filled vs. open circle (FODA says so
%                                        anyway), reinforced by blue vs. grey
%
% If a figure ever needs a third colour, the honest fix is to move that
% distinction onto one of those channels — not to add the colour here.
%
% ── USAGE ─────────────────────────────────────────────────────────────────
% % Shared colour palette for the ETHOS.TSAM manuscript figures.
%
% ── WHY THIS FILE EXISTS ──────────────────────────────────────────────────
% Reviewer request on the submitted manuscript:
%
%   "I would use as few different colors across your figures as possible. […]
%    I would make sure that the Jülich blue (RGB 2, 61, 107) is the leading
%    theme. That + a gray tone are the basic colors in all of your figures,
%    for shaded boxes e.g. at 70% transparency or so."
%
% Each figure used to declare its own \definecolor block, and they drifted:
% Fig. 1 carried two unrelated light blues plus a teal and an amber, Fig. 2 an
% amber and a warm amber tint. Declaring the palette ONCE and \input-ing it is
% what makes "the same colours in every figure" a property of the source
% rather than a thing to re-check by eye after every edit.
%
% ── THE PALETTE ───────────────────────────────────────────────────────────
% TWO base colours, and nothing else. Every fill in the figures is one of them
% mixed with white, which is exactly what "70 % transparency over a white
% page" means and what xcolor's `!nn` suffix computes: fzjbluetint is Jülich
% blue at 30 % opacity.
%
%   fzjblue  #023D6B  Jülich blue, RGB 2/61/107 — the leading theme. It draws
%                     the core: the tsam package itself, its internal edges,
%                     all primary text, and every mandatory feature.
%   fzjgrey  #3F4A55  the one grey — the periphery: external dependencies,
%                     optional or conditional elements, and legend chrome.
%
% ── THE CONSTRAINT THIS PUTS ON THE FIGURES ───────────────────────────────
% Both figures distinguish more than two things, and with only two hues the
% rest has to be carried by SHAPE, LINE STYLE or TYPOGRAPHY. That is a feature,
% not a compromise: all three survive a greyscale print, which a third hue
% would not. Concretely, what used to be colour is now
%
%   Fig. 1  object flow vs. dependency   solid vs. dashed (UML says so anyway)
%           pipeline phase vs. component round vs. sharp corners (UML Action)
%           optional stage               italic
%   Fig. 2  mandatory vs. optional       filled vs. open circle (FODA says so
%                                        anyway), reinforced by blue vs. grey
%
% If a figure ever needs a third colour, the honest fix is to move that
% distinction onto one of those channels — not to add the colour here.
%
% ── USAGE ─────────────────────────────────────────────────────────────────
% % Shared colour palette for the ETHOS.TSAM manuscript figures.
%
% ── WHY THIS FILE EXISTS ──────────────────────────────────────────────────
% Reviewer request on the submitted manuscript:
%
%   "I would use as few different colors across your figures as possible. […]
%    I would make sure that the Jülich blue (RGB 2, 61, 107) is the leading
%    theme. That + a gray tone are the basic colors in all of your figures,
%    for shaded boxes e.g. at 70% transparency or so."
%
% Each figure used to declare its own \definecolor block, and they drifted:
% Fig. 1 carried two unrelated light blues plus a teal and an amber, Fig. 2 an
% amber and a warm amber tint. Declaring the palette ONCE and \input-ing it is
% what makes "the same colours in every figure" a property of the source
% rather than a thing to re-check by eye after every edit.
%
% ── THE PALETTE ───────────────────────────────────────────────────────────
% TWO base colours, and nothing else. Every fill in the figures is one of them
% mixed with white, which is exactly what "70 % transparency over a white
% page" means and what xcolor's `!nn` suffix computes: fzjbluetint is Jülich
% blue at 30 % opacity.
%
%   fzjblue  #023D6B  Jülich blue, RGB 2/61/107 — the leading theme. It draws
%                     the core: the tsam package itself, its internal edges,
%                     all primary text, and every mandatory feature.
%   fzjgrey  #3F4A55  the one grey — the periphery: external dependencies,
%                     optional or conditional elements, and legend chrome.
%
% ── THE CONSTRAINT THIS PUTS ON THE FIGURES ───────────────────────────────
% Both figures distinguish more than two things, and with only two hues the
% rest has to be carried by SHAPE, LINE STYLE or TYPOGRAPHY. That is a feature,
% not a compromise: all three survive a greyscale print, which a third hue
% would not. Concretely, what used to be colour is now
%
%   Fig. 1  object flow vs. dependency   solid vs. dashed (UML says so anyway)
%           pipeline phase vs. component round vs. sharp corners (UML Action)
%           optional stage               italic
%   Fig. 2  mandatory vs. optional       filled vs. open circle (FODA says so
%                                        anyway), reinforced by blue vs. grey
%
% If a figure ever needs a third colour, the honest fix is to move that
% distinction onto one of those channels — not to add the colour here.
%
% ── USAGE ─────────────────────────────────────────────────────────────────
% \input{palette} AFTER \usepackage{tikz} (\colorlet needs xcolor, which tikz
% loads), and compile with the working directory set to this one — which is
% what every figure's "Regenerate:" note already does.

\definecolor{fzjblue}{HTML}{023D6B}
\definecolor{fzjgrey}{HTML}{3F4A55}

% The two shaded fills, named so that no figure ever spells a mix inline and
% the two cannot drift apart again.
\colorlet{fzjbluetint}{fzjblue!30}
\colorlet{fzjgreytint}{fzjgrey!20}
 AFTER \usepackage{tikz} (\colorlet needs xcolor, which tikz
% loads), and compile with the working directory set to this one — which is
% what every figure's "Regenerate:" note already does.

\definecolor{fzjblue}{HTML}{023D6B}
\definecolor{fzjgrey}{HTML}{3F4A55}

% The two shaded fills, named so that no figure ever spells a mix inline and
% the two cannot drift apart again.
\colorlet{fzjbluetint}{fzjblue!30}
\colorlet{fzjgreytint}{fzjgrey!20}
 AFTER \usepackage{tikz} (\colorlet needs xcolor, which tikz
% loads), and compile with the working directory set to this one — which is
% what every figure's "Regenerate:" note already does.

\definecolor{fzjblue}{HTML}{023D6B}
\definecolor{fzjgrey}{HTML}{3F4A55}

% The two shaded fills, named so that no figure ever spells a mix inline and
% the two cannot drift apart again.
\colorlet{fzjbluetint}{fzjblue!30}
\colorlet{fzjgreytint}{fzjgrey!20}
 AFTER \usepackage{tikz} (\colorlet needs xcolor, which tikz
% loads), and compile with the working directory set to this one — which is
% what every figure's "Regenerate:" note already does.

\definecolor{fzjblue}{HTML}{023D6B}
\definecolor{fzjgrey}{HTML}{3F4A55}

% The two shaded fills, named so that no figure ever spells a mix inline and
% the two cannot drift apart again.
\colorlet{fzjbluetint}{fzjblue!30}
\colorlet{fzjgreytint}{fzjgrey!20}
